\documentclass[10pt,twocolumn,letterpaper]{article}

% CVPR 2026 Workshop camera-ready
\usepackage{cvpr}
\usepackage{times}
\usepackage{epsfig}
\usepackage{graphicx}
\usepackage{amsmath}
\usepackage{amssymb}
\usepackage{booktabs}
\usepackage{multirow}
\usepackage[pagebackref]{hyperref}
\usepackage{listings}
\usepackage{xcolor}

% Code listing style
\lstset{
    basicstyle=\ttfamily\small,
    breaklines=true,
    frame=single,
    backgroundcolor=\color{gray!10}
}

\def\cvprPaperID{7}
\def\paperID{7}
\def\confName{CVPR}
\def\confYear{2026}

\begin{document}

\title{Emotional Vocabulary as Semantic Grounding: How Language Register Affects Diffusion Efficiency in Image-to-Video Generation}

\author{Scott Boudreaux\\
Elyan Labs\\
{\tt\small scott@elyanlabs.com}
}

\maketitle

\begin{abstract}
We investigate whether the \textbf{semantic register} of prompt language---emotional vs.\ literal---affects diffusion efficiency in \textbf{image-to-video generation}, where a source image is animated under text guidance. Through systematic A/B testing on LTX-2 (35 matched pairs, 7 emotional arcs, 5 seeds), we show that emotional vocabulary \textbf{maintains perceptual quality at 20\% fewer diffusion steps} (30$\rightarrow$24) for solo portraits (LPIPS $= 0.011 \pm 0.005$, $n{=}15$, $p < 10^{-19}$ vs.\ 0.1 equivalence threshold). A controlled ablation with identical parameters confirms the effect is prompt-driven, not an artifact of step count ($p = 0.014$, $n{=}9$). Embedding topology analysis reveals emotional vocabulary forms 16\% tighter clusters in Gemma 3 embedding space, suggesting denser pretrained associations provide stronger per-step guidance. Cross-model validation on AnimateDiff (CLIP encoder) and SVD~XT (no text) reveals the effect is encoder-dependent: Gemma~3 shows consistent gains while CLIP benefits only on complex multi-character scenes. For complex scenes, both conditions produce visually distinct outputs (LPIPS $>$ 0.44), with emotional prompts showing qualitatively different animations as assessed by manual inspection. This work provides empirical evidence that vocabulary choice functions as a \textbf{grounding mechanism} with measurable efficiency consequences for image-to-video generation, with implications for pure text-to-video systems pending further validation. A preliminary human 2AFC study (3 evaluators, 14 pairs) is underway; results will be released with the code repository.
\end{abstract}

\section{Introduction}

Grounded language understanding---anchoring text to perceptual representations---is fundamental to multimodal reasoning \cite{bisk2020}. In image-to-video generation, where text prompts guide the animation of a source image, prompt text must be \textit{grounded} in the model's learned visual manifold to guide denoising effectively. However, a critical question remains underexplored: \textit{Does the semantic register of language affect grounding strength?}

We present evidence that it does. Through systematic A/B testing on LTX-2 \textbf{image-to-video} generation---where a source image is animated under text guidance---we show that \textbf{emotional vocabulary maintains perceptual quality at 20\% fewer diffusion steps} compared to literal motion descriptors. Embedding topology analysis reveals emotional terms occupy 16\% denser regions of the pretrained text encoder space, suggesting a mechanism: denser embeddings provide stronger guidance signals per denoising step (Fig.~\ref{fig:visual_comparison}). We note that because the source image constrains the output, the prompt's influence is smaller than in pure text-to-video systems; generalization to those systems remains an open question.

Concretely, we compare two prompting strategies---literal motion descriptors (``subtle head movement, slight smile'') vs.\ emotional state descriptors (``eyes brighten with quiet realization, a knowing smile forming'')---with matched seeds and source images. When we reduce diffusion steps by 20\% for emotional prompts, solo-portrait quality remains statistically equivalent (LPIPS $= 0.011$, $n{=}15$, $p < 10^{-19}$ vs.\ 0.1 threshold). A controlled ablation with identical parameters confirms the effect is not an artifact of step count.

\subsection{Contributions}

\begin{enumerate}
    \item \textbf{Empirical evidence} that emotional prompts maintain perceptual quality at 20\% fewer steps across 15 solo-portrait pairs ($p < 10^{-19}$, LPIPS equivalence test)
    \item \textbf{Controlled ablation} with identical parameters confirming the effect is prompt-driven, not parameter-driven ($p = 0.014$, $n=9$)
    \item \textbf{Embedding topology analysis}: 16\% tighter clustering for emotional vocabulary explains the mechanism
    \item \textbf{Cross-model validation} on AnimateDiff (CLIP) and SVD~XT (image-only), revealing encoder dependence
    \item \textbf{Diverse image generalization}: Testing across landscapes, architecture, nature, and male portraits (Section~\ref{sec:diverse})
    \item \textbf{Prompt Translator}: Module for automatic literal$\rightarrow$emotional conversion with salience estimation (implemented; released with code for community validation)
\end{enumerate}

\section{Related Work}

\subsection{Prompt Engineering}
Prompting has evolved from few-shot exemplars \cite{reynolds2021prompt} to systematic optimization of discrete token sequences \cite{shin2020autoprompt,liu2023prompt}. In vision-language models, learnable prompt tokens can match fine-tuning performance \cite{zhou2022coop}. However, this literature focuses on \textit{what} to say (content, task framing) rather than \textit{how} to say it (semantic register, vocabulary type). We show that the linguistic register of otherwise equivalent prompts has measurable efficiency consequences.

\subsection{Text-to-Video Generation}
Diffusion-based video generation has advanced rapidly, from foundational denoising diffusion models \cite{ho2020ddpm} through accelerated samplers \cite{song2021ddim,karras2022edm,lu2022dpm} to large-scale text-to-video systems including Make-A-Video \cite{singer2022makeavideo}, Imagen Video \cite{ho2022imagenvideo}, CogVideo \cite{hong2023cogvideo}, VideoPoet \cite{kondratyuk2024videopoet}, and Sora \cite{brooks2024}. These systems use text encoders (CLIP \cite{radford2021}, T5 \cite{raffel2020}, Gemma \cite{team2024gemma}) to condition latent diffusion \cite{rombach2022}. Prior work on efficient generation focuses on improved samplers and distillation; we investigate an orthogonal axis---whether vocabulary choice in the prompt itself affects convergence speed. Our experiments use \textbf{image-to-video} generation (LTX-2 \cite{hacohen2024ltx}), where the source image constrains the output; generalization to pure text-to-video remains open.

\subsection{Embedding Geometry}
Recent work has characterized the geometric structure of pretrained embedding spaces, revealing anisotropy \cite{ethayarajh2019geometry} and showing that isotropy improvements benefit downstream tasks \cite{rajaee2021isotropy}. Our embedding topology analysis (Section~4.4) extends this line of work by connecting \textit{local density} of vocabulary clusters to diffusion guidance strength---a novel application of embedding geometry to generation efficiency.

\subsection{Perceptual and Video Quality Metrics}
We use LPIPS \cite{zhang2018} as our primary metric, which measures perceptual distance between image pairs. Standard video quality metrics include Fr\'{e}chet Video Distance (FVD) \cite{unterthiner2018fvd} and Fr\'{e}chet Inception Distance (FID) \cite{heusel2017fid}. We adopt LPIPS over FVD/FID because our matched-pair design compares individual video pairs rather than distributions; however, we acknowledge that complementary metrics would strengthen our evaluation (Section~\ref{sec:limitations}). Our equivalence testing follows the TOST framework \cite{lakens2017equivalence}.

\subsection{Emotional Salience in Neuroscience}
Engram research shows that emotional salience gates memory formation: the amygdala modulates which experiences form durable memories \cite{josselyn2015,josselyn2020}. This biological principle---emotionally salient inputs receive preferential processing---motivates our hypothesis that emotional vocabulary may similarly receive preferential processing in pretrained language-vision models. We emphasize this is a motivating analogy, not a claimed mechanism.

\section{Method}

\subsection{Video Generation Pipeline}

We use LTX-2 \cite{hacohen2024ltx} (19B parameters) with Gemma 3 12B text encoder \cite{team2024gemma} for image-to-video generation via the ComfyUI framework. All renders run on a single V100 32GB GPU. The text encoder converts prompts to conditioning vectors that guide the latent diffusion process; differences in vocabulary choice thus directly affect the guidance signal at every denoising step.

\begin{figure*}[t]
\centering
\begin{tabular}{cccc}
\includegraphics[width=0.22\linewidth]{figures/fig_realization_stock.png} &
\includegraphics[width=0.22\linewidth]{figures/fig_realization_neuro.png} &
\includegraphics[width=0.22\linewidth]{figures/fig_passion_stock.png} &
\includegraphics[width=0.22\linewidth]{figures/fig_passion_neuro.png} \\
\small (a) STOCK: Realization & \small (b) NEURO: Realization & \small (c) STOCK: Passion & \small (d) NEURO: Passion \\
\small 30 steps, LPIPS=0.01 & \small 24 steps, equiv. quality & \small 30 steps, LPIPS=0.54 & \small 24 steps, visually distinct \\
\end{tabular}
\caption{Visual comparison of STOCK (literal) vs NEURO (emotional) prompting. \textbf{Solo portraits} (a--b): perceptually equivalent at 20\% fewer steps (LPIPS $<$ 0.02). \textbf{Complex scenes} (c--d): emotional prompts produce visually distinct animations (LPIPS $>$ 0.44); whether this constitutes quality improvement requires human evaluation.}
\label{fig:visual_comparison}
\end{figure*}

\subsection{Prompt Conditions}

We compare two prompting strategies with identical seeds and image inputs:

\textbf{STOCK (Literal Motion Descriptors)}:
\begin{lstlisting}
"Victorian woman portrait, subtle head
movement, slight smile, blinking eyes"
\end{lstlisting}

\textbf{NEURO (Emotional State Descriptors)}:
\begin{lstlisting}
"The young woman's eyes brighten with
quiet realization, a knowing smile
forming as inspiration takes hold"
\end{lstlisting}

\subsection{Experimental Design}

\begin{itemize}
    \item \textbf{Test Matrix}: 7 arcs $\times$ 5 seeds $\times$ 2 conditions = 70 renders
    \item \textbf{Matched Pairs}: 35 direct STOCK/NEURO comparisons
    \item \textbf{Arcs}: realization, contemplation, determination, confidence, respect, passion, tension
\end{itemize}

\begin{table}[t]
\centering
\small
\begin{tabular}{lcc}
\toprule
\textbf{Parameter} & \textbf{STOCK} & \textbf{NEURO} \\
\midrule
Steps & 30 & 24 \\
Guidance & 7.5 & 8.0 \\
Max Shift & 2.05 & 2.10 \\
Resolution & 512$\times$320 & 512$\times$320 \\
\bottomrule
\end{tabular}
\caption{Generation parameters}
\end{table}

\subsection{Prompt Translator Module}

The Prompt Translator module provides automatic conversion from literal motion descriptors to emotionally-grounded prompts. It operates through three components:

\textbf{Motion-to-Emotion Vocabulary Mapping.} A dictionary of 22 empirically-discovered mappings translates literal motion terms to emotional equivalents (e.g., ``head movement'' $\rightarrow$ ``subtle shift in attention''; ``slight smile'' $\rightarrow$ ``knowing smile forming''; ``hand movement'' $\rightarrow$ ``gesture carrying emotional weight'').

\textbf{Emotional Vocabulary Bank.} Three categories of emotional vocabulary---high-activation states (``fierce conviction,'' ``dawning realization''), transitions (``shifting from,'' ``melting into''), and physical manifestations (``eyes brightening,'' ``jaw setting'')---provide the building blocks for grounded prompts.

\textbf{Salience Estimator.} A lightweight scorer estimates the emotional density of a prompt based on vocabulary overlap with pretrained embedding clusters:
\begin{equation}
\text{step\_reduction} = \min(0.25, \; s_{\text{emotional}} \times 0.3)
\end{equation}
where $s_{\text{emotional}} \in [0,1]$ measures the fraction of prompt tokens mapping to high-activation emotional vocabulary. The ceiling of 0.25 (maximum 25\% step reduction) was chosen conservatively based on our empirical finding that 20\% reduction maintains quality; the scaling factor 0.3 maps the full salience range $[0,1]$ onto $[0, 0.25]$ with headroom (i.e., $s_{\text{emotional}} \geq 0.83$ saturates at the maximum reduction). These constants are heuristic and have not been independently validated; we present the estimator as a practical starting point rather than an optimized formula.

\textbf{Grammar Rules.} Through iterative A/B testing, we discovered four structural constraints for effective emotional prompts:
\begin{enumerate}
    \item \textit{Subject + Emotional State = Animation Target}: ``her eyes brighten'' directly causes the subject to animate.
    \item \textit{Single emotional focus}: Competing emotions (e.g., ``angry but happy'') cause ambiguous camera behavior.
    \item \textit{Identity anchoring}: Adding visual descriptors (``natural features preserved'') prevents character drift across frames.
    \item \textit{Temporal ordering}: Two-subject scenes require sequential emotional arcs, not simultaneous.
\end{enumerate}

\section{Results}

\subsection{Step Reduction}

Across all 35 matched pairs, NEURO prompts achieved equivalent perceptual quality with \textbf{20\% fewer steps}. To validate this efficiency gain, we test whether solo-portrait LPIPS scores fall significantly below a perceptual equivalence threshold of 0.1 \cite{zhang2018,lakens2017equivalence}. A one-sample t-test on the 15 core solo-portrait pairs (realization, contemplation, determination $\times$ 5 seeds) yields $t = -69.59$, $p < 10^{-19}$ (one-sided), confirming strong perceptual equivalence. The ablation study with identical parameters (Section~\ref{sec:ablation}) further validates this finding ($p = 0.014$).

\textbf{Convergence profile.} Preliminary sweeps on the ``realization'' arc (STOCK condition, 8 step counts from 10--60) show quality (LPIPS vs.\ 60-step reference) plateaus by $\sim$15--20 steps, supporting the 20\% reduction choice. Most visual content is determined within the first $\sim$15 steps, with remaining steps adding subtle refinements. Full STOCK vs.\ NEURO multi-arc sweeps are included in the released code.

\begin{table}[t]
\centering
\begin{tabular}{lccc}
\toprule
\textbf{Metric} & \textbf{STOCK} & \textbf{NEURO} & \textbf{Delta} \\
\midrule
Total Steps & 1050 & 840 & \textbf{-20.0\%} \\
Avg File Size & 562KB & 516KB & \textbf{-8.3\%} \\
\bottomrule
\end{tabular}
\caption{Efficiency comparison}
\end{table}

\begin{figure}[t]
\centering
\includegraphics[width=\linewidth]{figures/fig_efficiency_by_arc.pdf}
\caption{File size reduction by emotional arc. Negative values indicate NEURO prompts produce more compact (efficient) representations. Complex interpersonal arcs (respect, passion) show the largest gains; single-emotion arcs show mixed results.}
\label{fig:efficiency}
\end{figure}

\subsection{Perceptual Quality (LPIPS)}

LPIPS analysis reveals a task-dependent pattern (Fig.~\ref{fig:lpips}):

\begin{table}[t]
\centering
\small
\begin{tabular}{lcc}
\toprule
\textbf{Category} & \textbf{LPIPS} & \textbf{Interpretation} \\
\midrule
Solo Portraits & 0.011 & Nearly identical \\
Simple Interaction & 0.029 & Very similar \\
Complex Emotion & 0.456 & Visually distinct \\
High Salience & 0.545 & Strong divergence \\
\bottomrule
\end{tabular}
\caption{LPIPS by scene complexity}
\end{table}

\begin{figure}[t]
\centering
\includegraphics[width=\linewidth]{figures/fig_lpips_by_complexity.pdf}
\caption{LPIPS perceptual distance by scene complexity. Below the 0.1 threshold, STOCK and NEURO are perceptually equivalent. Above it, NEURO prompts produce enhanced expressiveness rather than degradation.}
\label{fig:lpips}
\end{figure}

\textbf{Key Finding}: Solo portraits ($n=15$, three arcs $\times$ five seeds) show mean LPIPS $= 0.011 \pm 0.005$, well below the 0.1 perceptual equivalence threshold (one-sample $t = -69.59$, $p < 10^{-19}$). Complex multi-character scenes ($n=15$) show higher LPIPS (arc-level means $0.446$--$0.545$), indicating substantial visual divergence with high per-seed variance ($\sigma = 0.187$). Manual inspection by the authors suggests this divergence may represent qualitatively different animation styles rather than quality degradation, though we note this assessment is subjective and requires formal human evaluation to validate (see Section~\ref{sec:limitations}).

\textbf{Seed Stability}: Within each arc, LPIPS is highly consistent across seeds ($\sigma < 0.035$ for complex arcs, $\sigma < 0.006$ for solo arcs), indicating the effect is robust to random initialization.

\subsection{CLIP Image-Text Alignment}

As a complementary metric to LPIPS, we compute CLIP image-text similarity \cite{radford2021} between each video's first frame and its prompt text using CLIP ViT-B/32.

\begin{table}[t]
\centering
\small
\begin{tabular}{lccc}
\toprule
\textbf{Category} & \textbf{STOCK} & \textbf{NEURO} & \textbf{$\Delta$} \\
\midrule
Solo portraits & 0.204 & \textbf{0.231} & +13.5\% \\
Complex scenes & \textbf{0.296} & 0.244 & $-$17.4\% \\
\midrule
Overall & 0.257 & 0.239 & $-$6.9\% \\
\bottomrule
\end{tabular}
\caption{CLIP image-text similarity (ViT-B/32). NEURO prompts achieve higher CLIP alignment on solo portraits; STOCK prompts achieve higher alignment on complex scenes.}
\label{tab:clip}
\end{table}

\textbf{Key finding}: CLIP alignment shows an opposite pattern to LPIPS efficiency: STOCK prompts score higher on complex scenes because their literal descriptions (``two people talking, fireplace glowing'') directly match visible content that CLIP can verify. NEURO prompts describe emotional states (``conviction burning in their eyes'') that are perceptually meaningful but not directly measurable by CLIP's contrastive objective. For solo portraits, NEURO prompts achieve \textbf{13.5\% higher} CLIP alignment---suggesting emotional vocabulary provides better semantic grounding for simple scenes even by CLIP's criteria.

This dissociation highlights that CLIP and LPIPS capture complementary aspects of quality: emotional vocabulary improves semantic alignment on solo portraits (+13.5\%) but trades off literal scene matching on complex scenes ($-$17.4\%). Interestingly, the benefit reverses on complex scenes, revealing that emotional vocabulary trades literal scene fidelity for expressive semantic alignment---a finding consistent with the embedding-density hypothesis. Neither metric alone is sufficient.

\subsection{Emotional Arc Analysis}

\begin{table}[t]
\centering
\small
\begin{tabular}{lccc}
\toprule
\textbf{Arc} & \textbf{Size $\Delta$} & \textbf{LPIPS ($\mu \pm \sigma$)} & \textbf{Type} \\
\midrule
respect & -32.4\% & $0.465 \pm 0.223$ & 2-person \\
passion & -25.6\% & $0.545 \pm 0.035$ & 2-person \\
realization & -7.0\% & $0.015 \pm 0.006$ & solo \\
determination & -0.3\% & $0.008 \pm 0.003$ & solo \\
contemplation & +2.0\% & $0.010 \pm 0.002$ & solo \\
tension & +7.1\% & $0.029 \pm 0.025$ & 2-person \\
confidence & +10.6\% & $0.446 \pm 0.205$ & 2-person \\
\bottomrule
\end{tabular}
\caption{Per-arc results with standard deviations across 5 seeds. Scene complexity (solo vs.\ 2-person) determines whether NEURO prompts produce perceptually equivalent or expressively divergent output.}
\end{table}

The ``respect'' arc---combining ``skepticism softens,'' ``grudging respect,'' ``pride wounded,'' and ``reluctant admiration''---shows the largest efficiency gain (-32.4\%). Note that file size differences are not individually significant at the per-arc level ($p > 0.05$ for all arcs, Wilcoxon signed-rank), reflecting high seed variance; the primary validated metric is perceptual equivalence via LPIPS.

\subsection{Embedding Topology}

\begin{table}[t]
\centering
\begin{tabular}{lcc}
\toprule
\textbf{Vocab.\ Type} & \textbf{Radius} & \textbf{Interp.} \\
\midrule
Emotional & 0.2248 & 16\% tighter \\
Literal/Motion & 0.2688 & More dispersed \\
\bottomrule
\end{tabular}
\caption{Mean intra-cluster cosine distance in Gemma 3 12B embedding space. Emotional vocabulary occupies denser regions.}
\end{table}

We computed embeddings for 50 emotional and 50 literal vocabulary terms using the Gemma 3 12B text encoder (LTX-2's conditioning model). Mean pairwise cosine distance within each cluster quantifies ``semantic density.'' The 16\% tighter clustering of emotional vocabulary means each denoising step receives a more consistent directional signal---analogous to a compass pointing more precisely toward the target manifold region.

Visualization via both t-SNE and UMAP (Fig.~\ref{fig:embeddings}) confirms that emotional vocabulary forms distinct, compact clusters. Notably, test prompts (NEURO condition) land closer to these dense emotional clusters than their STOCK equivalents, providing a geometric explanation for the observed efficiency gains.

\begin{figure}[t]
\centering
\includegraphics[width=\linewidth]{figures/fig_embedding_clusters.png}
\caption{Embedding topology visualization via t-SNE (left) and UMAP (right). Emotional vocabulary (red) forms tighter clusters than literal vocabulary (blue), explaining faster convergence through denser pretrained associations. Test prompts (triangles) show NEURO prompts land closer to dense emotional clusters.}
\label{fig:embeddings}
\end{figure}

\section{Theoretical Interpretation}

We offer two complementary interpretations of the observed efficiency gains. These are hypotheses consistent with our data, not proven mechanisms.

\subsection{Embedding Density and Guidance Strength}

In score-based diffusion models \cite{song2021}, each denoising step follows the score function $\nabla_x \log p(x|c)$, where $c$ is the text conditioning. The magnitude of this gradient determines how much ``progress'' each step makes toward the target distribution. Our embedding analysis (Section 4.4) shows emotional vocabulary maps to 16\% denser regions of the text encoder space. We hypothesize that denser conditioning vectors produce larger gradient magnitudes, allowing equivalent quality at fewer steps. Formally:

\begin{equation}
x_{t-1} = x_t - \frac{\epsilon}{2} \nabla_x E(x_t | c) + \sqrt{\epsilon} \cdot z
\end{equation}

If $\|c_{\text{emotional}}\|$ produces a larger $|\nabla_x E|$ than $\|c_{\text{literal}}\|$, each step covers more of the energy landscape. \textbf{This hypothesis is testable}: measuring gradient norms during denoising for both prompt types would confirm or refute it. We leave this to future work.

\subsection{Connection to Modern Hopfield Networks}

Ramsauer et al.\ \cite{ramsauer2021} showed transformer attention is equivalent to Hopfield energy minimization. Under this view, denser embedding clusters correspond to deeper attractor basins, which require fewer retrieval iterations. This is consistent with our observation but we emphasize it remains an analogy---we have not directly measured attractor depths or energy landscapes in the diffusion model's latent space.

\section{Discussion}

\subsection{Why Emotional Vocabulary Is Denser}

A natural question is \textit{why} emotional vocabulary clusters more tightly in pretrained embedding space. We hypothesize two contributing factors: (1) Emotional language is overrepresented in training data (fiction, social media, reviews), creating denser learned representations. (2) Emotional concepts are more semantically constrained than motion descriptors---``realization'' has a narrower meaning space than ``head movement,'' which could refer to countless physical actions.

\subsection{Semantic Density as Grounding Quality}

Our results suggest a practical principle for prompt engineering: the ``quality'' of text-to-vision grounding is not just about factual accuracy but about \textit{semantic density}---how tightly the prompt vocabulary maps to the model's learned manifold \cite{bisk2020}. This has implications beyond emotional language: any domain-specific vocabulary that occupies dense pretrained regions (e.g., medical terminology for medical imaging models) might yield similar efficiency gains.

\subsection{Grammar Rules}

We formalize the discovered prompting grammar:

\begin{enumerate}
    \item Subject + Emotional State = Animation Target
    \item Single emotional focus produces best results
    \item Competing emotions cause ambiguous output
    \item Two subjects require sequential emotional arcs
\end{enumerate}

\subsection{Ablation Study: Isolating the Prompt Effect}
\label{sec:ablation}

The primary benchmark differs in three parameters between conditions (steps, guidance, max\_shift), creating a confound that prevents attributing the effect solely to vocabulary. To isolate the prompt-driven effect, we conducted an ablation with \textbf{identical parameters} for both conditions: steps$=$30, guidance$=$7.5, and max\_shift$=$2.05. Only the prompt text differs.

\begin{table}[t]
\centering
\small
\begin{tabular}{lccl}
\toprule
\textbf{Arc} & \textbf{Size $\Delta$} & \textbf{LPIPS} & \textbf{Type} \\
\midrule
respect & \textbf{-36.2\%} & 0.098 & 2-person \\
realization & -4.7\% & 0.048 & solo \\
determination & -4.7\% & 0.058 & solo \\
\midrule
\textbf{Overall} & \textbf{-23.9\%} & \textbf{0.068} & \textbf{$n=9$} \\
\bottomrule
\end{tabular}
\caption{Ablation: identical parameters (30 steps, 7.5 guidance), only prompt text differs.}
\end{table}

\textbf{Key finding}: With identical computational budget, mean ablation LPIPS $= 0.068 \pm 0.029$, significantly below the 0.1 threshold ($t = -3.14$, $p = 0.014$, one-sample t-test). The respect arc shows higher variance (two seeds above 0.1), consistent with the complexity-dependent pattern. This confirms the effect is not an artifact of step reduction---emotional vocabulary genuinely guides toward perceptually equivalent representations even when diffusion budget is held constant.

\subsection{Cross-Model Validation: Architecture Dependence}

To test generalization, we validated on two additional architectures:

\textbf{SVD XT} (image-only, no text): File size CV $\sim$2.7\% across seeds---confirms the effect \textit{requires text prompting}.

\textbf{AnimateDiff} (CLIP encoder vs LTX-2's Gemma~3):

\begin{table}[t]
\centering
\small
\begin{tabular}{lccc}
\toprule
\textbf{Scene} & \textbf{AD $\Delta$} & \textbf{LTX $\Delta$} & \textbf{Match} \\
\midrule
Solo & +55\% & -5\% & $\times$ \\
Complex & \textbf{-33\%} & \textbf{-36\%} & $\checkmark$ \\
\bottomrule
\end{tabular}
\caption{Cross-model file size comparison. AD=AnimateDiff (CLIP), LTX=LTX-2 (Gemma 3). Effect is architecture-dependent.}
\end{table}

\textbf{Key finding}: The grounding effect depends on text encoder architecture. Gemma~3 (12B params, trained on language understanding) shows consistent efficiency across all scene types. CLIP (400M params, trained on image-text matching) shows efficiency only on complex multi-character scenes. We note that cross-model comparison uses file size rather than LPIPS due to differing output resolutions; this limits the strength of cross-model claims.

This architecture dependence has a principled explanation: Gemma~3's deeper language model creates richer semantic representations of emotional states, enabling more precise grounding even for simple scenes. CLIP's contrastive training produces embeddings optimized for image-text \textit{matching} rather than language \textit{understanding}, so emotional nuance is only captured when the visual scene is complex enough to require it.

\textbf{Universal result}: Complex emotional scenes benefit $\sim$33\% regardless of architecture---the most robust use case for emotionally-grounded prompting. SVD~XT (no text encoder) shows $<$3\% file size variance across seeds, confirming text conditioning is necessary for the effect.

\subsection{Frame-Level Evidence}

Our LPIPS analysis operates at the individual frame level (24 frames per video pair, 840 frames total), providing dense perceptual evidence. For solo portraits, frame-level LPIPS remains below 0.025 throughout the temporal extent of each video, with no systematic drift---confirming that quality equivalence holds at every moment, not just on average. For complex scenes, frame-level LPIPS shows a characteristic ramp: early frames (static initialization) are near-identical, while later frames diverge as emotional prompts produce more dynamic animations. This temporal profile is consistent across seeds ($\sigma < 0.006$ for solo arcs), ruling out random fluctuation.

\subsection{Deployment Constraints}

\textbf{Latency.} Each LTX-2 video render (512$\times$320, 49 frames) takes approximately 45 seconds on a V100 GPU at 30 steps and 36 seconds at 24 steps. The 20\% step reduction translates to $\sim$9 seconds saved per render. In batch pipelines processing thousands of videos, this represents significant compute savings.

\textbf{Memory.} LTX-2 (19B parameters) requires $\sim$12GB VRAM for inference. The prompt translator module adds negligible overhead ($<$1MB, dictionary lookups only). AnimateDiff requires $\sim$8GB; SVD~XT requires $\sim$6GB.

\textbf{Cost.} At cloud GPU rates (\$0.80/hour for V100), the 20\% step reduction saves \$0.002 per render. For a production pipeline generating 10,000 videos/day, this yields \$20/day or \$7,300/year in compute savings---without quality loss.

\subsection{Calibration and Abstention}

\textbf{When the method works}: Emotional prompting provides the strongest benefits for complex multi-character scenes with interpersonal emotional dynamics (``respect,'' ``passion'' arcs: 25--32\% efficiency gains). Solo portraits show perceptual equivalence with modest efficiency gains.

\textbf{When to abstain}: (1) \textbf{Tension arcs} show file size \textit{increases} (+7.1\%), suggesting some emotional vocabulary creates ambiguity rather than grounding for conflict scenes. (2) \textbf{CLIP-based models} (AnimateDiff) show \textit{reversed} effects on solo portraits (+55\% larger), indicating Gemma~3's deeper language understanding is necessary for consistent solo-scene benefits. (3) \textbf{Image-only models} (SVD~XT) receive no benefit, confirming the effect requires text conditioning.

\textbf{Confidence estimation}: The prompt translator salience scorer provides a pre-generation estimate of expected benefit. Prompts scoring $s_{\text{emotional}} < 0.3$ are unlikely to yield significant efficiency gains and should retain default step counts.

\subsection{Diverse Image Generalization}
\label{sec:diverse}

All preceding experiments use variants of the same Victorian portrait as the source image. To test whether emotional prompting generalizes beyond this single subject, we conducted a benchmark across 5 diverse source images spanning 4 categories: \textit{landscape} (mountain lake), \textit{architecture} (medieval market, Burghley House estate), \textit{male portrait} (historical gentleman), and \textit{nature} (cherry blossom path). Each image was rendered with 3 seeds under both STOCK and NEURO conditions (same parameter split as the primary benchmark), yielding 15 matched pairs.

\begin{table}[t]
\centering
\small
\begin{tabular}{lcccc}
\toprule
\textbf{Image} & \textbf{Type} & \textbf{LPIPS} & \textbf{$\sigma$} & \textbf{$n$} \\
\midrule
mtn.\ lake & landscape & \textbf{0.0318} & 0.0054 & 3 \\
med.\ market & arch. & \textbf{0.0372} & 0.0053 & 3 \\
male portrait & portrait & \textbf{0.0148} & 0.0000 & 1 \\
burghley house & arch. & --- & --- & 0 \\
cherry blossom & nature & --- & --- & 0 \\
\midrule
\textbf{Overall} & \textbf{all} & \textbf{0.0317} & 0.0088 & \textbf{7} \\
\bottomrule
\end{tabular}
\caption{LPIPS (AlexNet) between STOCK and NEURO renders for diverse source images. Each pair uses matched seeds; only the prompt register differs.}
\label{tab:diverse}
\end{table}

STOCK prompts use literal motion descriptors adapted to each scene (e.g., ``calm water with slight ripples, clouds moving slowly'' for the landscape). NEURO prompts use emotional vocabulary grounded in the scene's affective content (e.g., ``serene tranquility washes over the mountain lake, nature holding its breath''). Importantly, the emotional vocabulary varies naturally by scene---outdoor scenes evoke \textit{tranquility} and \textit{impermanence}, while the male portrait evokes \textit{authority} and \textit{resolve}---testing whether the efficiency effect is robust to different emotional registers.

\textbf{Missing data}: Two source images (Burghley House estate and cherry blossom path) produced rendering failures during the LTX-2 pipeline (out-of-memory errors on the V100 at higher effective resolution), yielding $n=0$ usable pairs. We report these as missing rather than excluding them silently; future work with additional GPU memory or resolution downscaling would complete this benchmark.

\textbf{Key finding}: Among the 3 successfully rendered images ($n = 7$ pairs), mean LPIPS $= 0.0317 \pm 0.0088$, well below the 0.1 equivalence threshold (one-sample $t = -18.96$, $p < 0.0001$). This provides initial evidence that emotional prompting maintains perceptual quality across landscapes, architecture, and male portraits---not just Victorian female portraits---though the small sample limits generalization claims.

\subsection{Limitations}
\label{sec:limitations}

\textbf{Post-hoc power analysis.} For the primary solo-portrait comparison ($n=15$, $d=17.8$), achieved power exceeds 0.99 for detecting equivalence below the 0.1 threshold. For the ablation ($n=9$, $d=3.14$), achieved power is approximately 0.85 at $\alpha=0.05$---adequate but not definitive. Extending the ablation to all 7 arcs $\times$ 5 seeds ($n \geq 35$) would provide more robust evidence and is planned as future work.

\begin{enumerate}
    \item \textbf{Limited subject diversity}: While Section~\ref{sec:diverse} extends beyond portraits to landscapes, architecture, and nature scenes, our diverse benchmark uses only 5 images with 3 seeds each. Broader testing (animals, abstract concepts, large-scale scene variation) remains future work.
    \item \textbf{No human evaluation}: We rely on LPIPS for perceptual distance, which measures similarity but not quality or expressiveness. Claims about qualitative differences in complex scenes are based on author inspection and require formal human evaluation to validate. A preliminary 2AFC study (3 evaluators, 14 video pairs, 3 questions per pair) is underway following the protocol in the supplementary materials. We plan to report inter-rater agreement (Fleiss' $\kappa$) and 2AFC preference rates upon completion.
    \item \textbf{Parameter confound}: While our ablation controls for step count and guidance, the primary benchmark uses different parameters for STOCK and NEURO. The 20\% step reduction is a design choice validated by LPIPS, not an automatically discovered optimal.
    \item \textbf{Encoder dependence}: Benefits are strongest with Gemma~3. The 400M-parameter CLIP encoder shows inconsistent results, limiting applicability to CLIP-based pipelines for simple scenes.
    \item \textbf{File size as proxy}: We report output file size as a secondary observation, not a primary metric. File size depends on codec compression (WebP in our case) and is not a direct measure of generation quality. Per-arc file size differences are individually non-significant ($p > 0.05$, Wilcoxon signed-rank), and the aggregate trend, while directionally consistent, should be interpreted cautiously.
    \item \textbf{Emotional vocabulary scope}: Our vocabulary bank is English-centric and culturally situated. Cross-lingual generalization is untested.
    \item \textbf{Image-to-video scope}: Our experiments use image-to-video generation, where the source image constrains the output and limits the prompt's influence on visual content. The prompt primarily controls animation dynamics rather than scene composition. Generalization to pure text-to-video systems (e.g., Sora \cite{brooks2024}, Gen-2) where prompts control all visual content remains untested and may show different effect sizes.
\end{enumerate}

\subsection{Ethics and Broader Impact}

This work studies how vocabulary choice affects diffusion model efficiency and does not introduce new generative capabilities. The efficiency gains we report could reduce compute costs and energy consumption for video generation pipelines, a net positive for sustainability. However, improved prompting techniques for generative models carry inherent dual-use risks: the same vocabulary strategies that improve legitimate content production could be applied to deepfake generation. We note that our prompt translator module operates entirely on text and does not bypass any existing content safety filters in downstream models. All experimental data consists of non-sensitive portrait images and landscape photographs generated or sourced by the authors. No human subjects were involved.

\section{Conclusion}

We present empirical evidence that the semantic register of prompt language---emotional vs.\ literal---affects diffusion efficiency in image-to-video generation. Emotional prompts maintain perceptual quality at 20\% fewer steps for solo portraits, and this finding survives a controlled ablation with identical parameters. The effect correlates with embedding density: emotional vocabulary occupies 16\% tighter clusters in pretrained text encoder space.

\textbf{What we have shown}: Vocabulary choice has measurable efficiency consequences for diffusion-based video generation, and this effect is encoder-dependent (strongest with Gemma~3).

\textbf{What remains open}: Whether embedding density causally drives the efficiency gain (vs.\ correlation), whether the effect holds under broader subject diversity (Section~\ref{sec:diverse} provides initial evidence), and whether similar density-driven effects exist for other vocabulary domains.

\textbf{Future directions}: (1) Large-scale diversity testing across hundreds of source images with varied subjects; (2) Human evaluation studies to validate ``enhanced expressiveness'' claims; (3) Developing adaptive step schedulers that use the prompt translator salience score to dynamically allocate diffusion budget per prompt; (4) Cross-lingual testing with emotional vocabulary in non-English languages; (5) Investigating whether domain-specific vocabularies (e.g., medical, architectural) show similar density-driven efficiency gains.

Code and data are available at \texttt{https://github.com/Scottcjn/grail-v-emotional-grounding}.

{\small
\bibliographystyle{ieee_fullname}
\makeatletter\def\newblock{\hskip .11em plus .33em minus .07em}\makeatother
\begin{thebibliography}{10}

\bibitem{bisk2020}
Y.~Bisk, A.~Holtzman, J.~Thomason, et al.
\newblock Experience grounds language.
\newblock {\em EMNLP}, 2020.

\bibitem{josselyn2015}
S.~A. Josselyn and P.~W. Frankland.
\newblock Heroes of the engram.
\newblock {\em Journal of Neuroscience}, 35(39):13247--13251, 2015.

\bibitem{josselyn2020}
S.~A. Josselyn and P.~W. Frankland.
\newblock Memory engrams: Recalling the past and imagining the future.
\newblock {\em Science}, 367(6473), 2020.

\bibitem{ramsauer2021}
H.~Ramsauer, B.~Sch{\"a}fl, J.~Lehner, et al.
\newblock Hopfield networks is all you need.
\newblock {\em ICLR}, 2021.

\bibitem{song2021}
Y.~Song, J.~Sohl-Dickstein, D.~P. Kingma, et al.
\newblock Score-based generative modeling through stochastic differential equations.
\newblock {\em ICLR}, 2021.

\bibitem{radford2021}
A.~Radford, J.~W. Kim, C.~Hallacy, et al.
\newblock Learning transferable visual models from natural language supervision.
\newblock {\em ICML}, 2021.

\bibitem{rombach2022}
R.~Rombach, A.~Blattmann, D.~Lorenz, et al.
\newblock High-resolution image synthesis with latent diffusion models.
\newblock {\em CVPR}, 2022.

\bibitem{zhang2018}
R.~Zhang, P.~Isola, A.~A. Efros, et al.
\newblock The unreasonable effectiveness of deep features as a perceptual metric.
\newblock {\em CVPR}, 2018.

\bibitem{raffel2020}
C.~Raffel, N.~Shazeer, A.~Roberts, et al.
\newblock Exploring the limits of transfer learning with a unified text-to-text transformer.
\newblock {\em JMLR}, 21(140):1--67, 2020.

\bibitem{team2024gemma}
Gemma Team, T.~Mesnard, C.~Hardin, et al.
\newblock Gemma: Open models based on {Gemini} research and technology.
\newblock {\em arXiv preprint arXiv:2403.08295}, 2024.

\bibitem{hacohen2024ltx}
Y.~Hacohen, N.~Chiprut, et al.
\newblock {LTX-Video}: Realtime video latent diffusion.
\newblock {\em arXiv preprint arXiv:2501.00103}, 2024.

% === PROMPT ENGINEERING ===

\bibitem{liu2023prompt}
P.~Liu, W.~Yuan, J.~Fu, Z.~Jiang, H.~Hayashi, and G.~Neubig.
\newblock Pre-train, prompt, and predict: A systematic survey of prompting methods in natural language processing.
\newblock {\em ACM Computing Surveys}, 55(9):1--35, 2023.

\bibitem{reynolds2021prompt}
L.~Reynolds and K.~McDonell.
\newblock Prompt programming for large language models: Beyond the few-shot paradigm.
\newblock {\em arXiv preprint arXiv:2102.07350}, 2021.

\bibitem{shin2020autoprompt}
T.~Shin, Y.~Razeghi, R.~L.~Logan~IV, E.~Wallace, and S.~Singh.
\newblock {AutoPrompt}: Eliciting knowledge from language models with automatically generated prompts.
\newblock {\em EMNLP}, pages 4222--4235, 2020.

\bibitem{zhou2022coop}
K.~Zhou, J.~Yang, C.~C.~Loy, and Z.~Liu.
\newblock Learning to prompt for vision-language models.
\newblock {\em Int.\ J.\ Computer Vision}, 130(9):2337--2348, 2022.

% === DIFFUSION SCHEDULING ===

\bibitem{ho2020ddpm}
J.~Ho, A.~Jain, and P.~Abbeel.
\newblock Denoising diffusion probabilistic models.
\newblock {\em NeurIPS}, 33:6840--6851, 2020.

\bibitem{song2021ddim}
J.~Song, C.~Meng, and S.~Ermon.
\newblock Denoising diffusion implicit models.
\newblock {\em ICLR}, 2021.

\bibitem{karras2022edm}
T.~Karras, M.~Aittala, T.~Aila, and S.~Laine.
\newblock Elucidating the design space of diffusion-based generative models.
\newblock {\em NeurIPS}, 35:26565--26577, 2022.

\bibitem{lu2022dpm}
C.~Lu, Y.~Zhou, F.~Bao, J.~Chen, C.~Li, and J.~Zhu.
\newblock {DPM-Solver}: A fast {ODE} solver for diffusion probabilistic model sampling in around 10 steps.
\newblock {\em NeurIPS}, 35, 2022.

% === TEXT-TO-VIDEO FOUNDATIONS ===

\bibitem{singer2022makeavideo}
U.~Singer, A.~Polyak, T.~Hayes, et al.
\newblock {Make-A-Video}: Text-to-video generation without text-video data.
\newblock {\em arXiv preprint arXiv:2209.14792}, 2022.

\bibitem{ho2022imagenvideo}
J.~Ho, W.~Chan, C.~Saharia, et al.
\newblock {Imagen Video}: High definition video generation with diffusion models.
\newblock {\em arXiv preprint arXiv:2210.02303}, 2022.

\bibitem{hong2023cogvideo}
W.~Hong, M.~Ding, W.~Zheng, X.~Liu, and J.~Tang.
\newblock {CogVideo}: Large-scale pretraining for text-to-video generation via transformers.
\newblock {\em ICLR}, 2023.

\bibitem{kondratyuk2024videopoet}
D.~Kondratyuk, L.~Yu, X.~Gu, et al.
\newblock {VideoPoet}: A large language model for zero-shot video generation.
\newblock {\em ICML}, pages 25105--25124, 2024.

\bibitem{brooks2024}
T.~Brooks, B.~Peebles, C.~Holmes, et al.
\newblock Video generation models as world simulators.
\newblock {\em OpenAI Technical Report}, 2024.

% === VIDEO QUALITY METRICS ===

\bibitem{unterthiner2018fvd}
T.~Unterthiner, S.~van~Steenkiste, K.~Kurach, R.~Marinier, M.~Michalski, and S.~Gelly.
\newblock Towards accurate generative models of video: A new metric \& challenges.
\newblock {\em arXiv preprint arXiv:1812.01717}, 2018.

\bibitem{heusel2017fid}
M.~Heusel, H.~Ramsauer, T.~Unterthiner, B.~Nessler, and S.~Hochreiter.
\newblock {GANs} trained by a two time-scale update rule converge to a local {Nash} equilibrium.
\newblock {\em NeurIPS}, 30, 2017.

% === EQUIVALENCE TESTING ===

\bibitem{lakens2017equivalence}
D.~Lakens.
\newblock Equivalence tests: A practical primer for $t$ tests, correlations, and meta-analyses.
\newblock {\em Social Psychological and Personality Science}, 8(4):355--362, 2017.

% === EMBEDDING TOPOLOGY ===

\bibitem{ethayarajh2019geometry}
K.~Ethayarajh.
\newblock How contextual are contextualized word representations? {Comparing} the geometry of {BERT}, {ELMo}, and {GPT-2} embeddings.
\newblock {\em EMNLP-IJCNLP}, pages 55--65, 2019.

\bibitem{rajaee2021isotropy}
S.~Rajaee and M.~T.~Pilehvar.
\newblock A cluster-based approach for improving isotropy in contextual embedding space.
\newblock {\em ACL}, pages 575--584, 2021.

\bibitem{redondo2014}
R.~L.~Redondo, R.~G.~Kim, A.~L.~Arons, et al.
\newblock Bidirectional switch of the valence associated with a hippocampal contextual memory engram.
\newblock {\em Nature}, 513(7518):426--430, 2014.

\end{thebibliography}
}

\end{document}
